\tocsection{\textbf{10/28/26 Lecture 26:} Functions of Bounded Variation}

\textbf{Recall:} If $F : \R \to \R$ is increasing, $F$ is differentiable a.e.\\

For $F : \R \to \R$, consider
\[ T_F(x) = \sup\left\{\sum_{j = 1}^{n} \abs{F(x_j) - F(x_{j - 1})} \col n \in \N,\ x_0 < x_1 < \cdots < x_n = x\right\} \]
and if $T_F(\infty) < \infty$, $F \in \tt{BV}$.\\

Now, if $x_{j - 1} < x' < x_j$,
\[ \abs{F(x_j) - F(x_{j - 1})} \leq \abs{F(x_j) - F(x')} + \abs{F(x') - F(x_{j - 1})} \]
and if $a < b$,
\[ T_F(b) = T_F(a) + \sup\left\{\sum \abs{F(x_j) - F(x_{j - 1})} \col a = x_0 < \cdots < x_n = b\right\} \]

\begin{definition}[\textbf{Definition 26.1:} Bounded Variation on an Interval]
    We say $F$ is of \emph{bounded variation on $[a, b]$} if
    \[ T_F(b) - T_F(a) < \infty \]
    and we write $F \in \tt{BV}([a, b])$.
\end{definition}

\begin{lemma}[\textbf{Proposition 26.2}]
    If $F \in \tt{BV}$, then $T_F + F$ and $T_F - F$ are increasing.
\end{lemma}
\begin{proof}{Proof:}
    Fix $x < y$. By definition, choose $x_0 < \cdots < x_n = x$ so that $\sum_{j} \abs{F(x_j) - F(x_{j - 1})} \geq T_F(x) - \epsilon$. Then,
    \begin{align*}
        T_F(y) + F(y) &\geq \sum_{j = 1}^{n} \abs{F(x_j) - F(x_{j - 1})} + \abs{F(y) - F(x)} + F(y)\\
        &\geq T_F(x) - \epsilon + F(x)
    \end{align*}
    so taking $\epsilon \to 0$, we get what we want. A similar argument works for $T_{F} - F$. 

    $\qed$
\end{proof}

\begin{theorem}[\textbf{Theorem 26.3:} Classification of BV]
    A function $F : \R \to \R$ is BV iff it may be written as a difference of bounded increasing functions. In particular, $F \in \tt{BV} \implies F$ is differentiable almost everywhere and $F$ is bounded.
\end{theorem}
\begin{proof}{Proof:}
    $(\Rightarrow)$ By the proposition, we can write
    \[ F = \frac{1}{2}(T_F + F) - \frac{1}{2}(T_F - F) \]
    and both are increasing. To show boundedness, we see that, for $x < y$,

    \[ F(x) - F(y) \leq T_F(y) - T_F(x) \leq T_F(\infty) - T_F(-\infty) < \infty \]
    and similarly
    \[ F(y) - F(x) \leq T_F(y) - T_F(x) \leq T_F(\infty) - T_F(-\infty) < \infty \]
    so we see $F$ is bounded, so $T_F + F$ and $T_F - F$ are bounded.\\


    $(\Leftarrow)$ Bounded increasing functions have bounded variation, and as $T_{H - G} \leq T_{H} + T_{G}$ by the triangle inequality, we see that the difference also has bounded variation.

    $\qed$
\end{proof}

Recall if $F$ is increasing and right-continuous, we get the LS measure $m_F$ by $m_F((a, b]) = F(b) - F(a)$.

\begin{lemma}[\textbf{Proposition 26.4}]
    Let $m$ be a positive Borel measure with the property that $m(A) < \infty$ if $A$ is bounded. Then, if
    \[ F(x) = \left\{\begin{matrix} m((0, x]) & x > 0 \\ 0 & x = 0 \\ -m((x, 0]) & x < 0 \end{matrix}\right. \]
    then $F$ is increasing, right continuous, and $m = m_F$.
\end{lemma}
\begin{proof}{Proof:}

    Increasing follows from $A \subseteq B \implies m(A) \leq m(B)$, and right continuity comes from
    \[ m((0, x]) = \lim_{h \to 0^{+}} m((0, x + h]) \]
    It is then clear that $m = m_F$.

    $\qed$
\end{proof}

\begin{lemma}[\textbf{Corollary 26.5}]
    If $m$ is a Borel measure with $m(A) < \infty$ for bounded $A$, then $m$ is regular!
\end{lemma}

This gives us a correspondence $m_F \longleftrightarrow F$, but not quite, as if $G = F + C$ for a constant $C$, then $m_F = m_G$.\\

To fix this, we will normalize as follows: let
\[ \tt{NBV} = \{F \in \tt{BV} \col F \tt{ is right continuous and } F(-\infty) = 0\} \]

\begin{theorem}[\textbf{Theorem 26.6:} NBV Functions and Signed Borel Measures]
    If $\mu$ is a finite signed measure on $\R$, then $F(x) = \mu((-\infty, x])$ is NBV. Conversely, if $F \in \tt{NBV}$, then there is a unique signed Borel measure $\mu$ with $\mu((a, b]) = F(b) - F(a)$.
\end{theorem}

Before we prove this, we need a lemma.

\begin{lemma}[\textbf{Lemma 26.7}]
    If $F \in \tt{BV}$ and is right-continuous, then so is $T_F(x)$. \
\end{lemma}
\begin{proof}{Proof:}
    Fix $x$ and let $\alpha = T_F(x^{+}) - T_F(x)$. We want to show $\alpha = 0$. Given $\epsilon > 0$, choose $\delta$ so that
    \[ \abs{F(x + h) - F(x)} < \epsilon \quand T_F(x + h) - T_F(x^{+}) < \epsilon \]
    for $0 < h < \delta$ by right-continuity and the definition of limits. Find a partition $x = x_0 < \cdots < x_n = x + h$ so that
    \begin{align*}
        \sum_{j = 1}^{n} \abs{F(x_j) - F(x_{j - 1})} &\geq \frac{3}{4}\abs{T_F(x + h) - T_F(x)}\\
        &\geq \frac{3}{4}\alpha
    \end{align*}
    \[ \sum_{j = 2}^{n} \abs{F(x_j) - F(x_{j - 1})} \geq \frac{3}{4}\alpha - \abs{F(x_1) - F(x)} \geq \frac{3}{4}\alpha - \epsilon \]

    Now, find $x = y_0 < \cdots < y_m = x_1$ so that
    \[ \sum_{j = 1}^{m} \abs{F(y_j) - F(y_{j - 1})} \geq \frac{3}{4}\abs{T_F(x_1) - T_F(x)} \geq \frac{3}{4}\alpha \]
    Thus,
    \[ \sum_{j = 1}^{m} \abs{F(y_j) - F(y_{j - 1})} + \sum_{j = 2}^{n} \abs{F(x_j) - F(x_{j - 1})} \geq \frac{3}{2}\alpha - \epsilon \]
    But as this is just a partition (of $[x, x + h]$),
    \begin{align*}
        \sum_{j = 1}^{m} \abs{F(y_j) - F(y_{j - 1})} + \sum_{j = 2}^{n} \abs{F(x_j) - F(x_{j - 1})} &\leq T_F(x + h) - T_F(x)\\
        &= (T_F(x + h) - T_F(x^{+})) + (T_F(x^{+}) - T_F(x))\\
        &< \epsilon + \alpha
    \end{align*}
    so
    \[ \frac{3}{2}\alpha - \epsilon \leq T_F(x + h) - T_F(x) < \epsilon + \alpha \]
    so taking $\epsilon \to 0$, we must see $\alpha = 0$. Since $T_F(x)$ is increasing and bounded, $T_F \in \tt{BV}$.

    $\qed$
\end{proof}

\begin{proof}{Proof of Theorem 26.6:}
    $(\Rightarrow)$ If $\mu$ is a finite signed Borel measure, write $\mu = \mu^{+} - \mu^{-}$ and let
    \[ F^{\pm}(x) = \mu^{\pm}((-\infty, x]) \]
    Each of these are BV and $F^{\pm}(-\infty) = 0$ by continuity. They are also right continuous by continuity.\\

    $(\Leftarrow)$ Write $F = F^{+} - F^{-} = \frac{1}{2}(T_F + F) - \frac{1}{2}(T_F - F)$. $F^{\pm}$ are both increasing and right continuous from above, so we get the LS measures $m_{F^{+}}$ and $m_{F^{-}}$, and it can be checked that
    \[ m = m_{F^{+}} - m_{F^{-}} \]
    has the property we want.

    $\qed$
\end{proof}
